\documentclass[11pt]{article}

\usepackage[margin=1.1in]{geometry}
\usepackage{amsmath}
\usepackage{amssymb}
\usepackage{amsthm}
\usepackage{booktabs}
\usepackage{microtype}
\usepackage[round,authoryear]{natbib}
\usepackage[colorlinks=true,allcolors=blue]{hyperref}

\newtheoremstyle{itheads}{\topsep}{\topsep}%
  {\itshape}{}{\itshape}{.}{ }{\thmname{#1}\thmnumber{ #2}%
  \thmnote{ (#3)}}
\theoremstyle{itheads}
\newtheorem{proposition}{Proposition}

\title{\textbf{A Posterior Predictive for Pass@$k$}\\[0.4em]
\large Established Empirical Bayes, Measured on Released Rollouts}
\author{Peter Cotton}
\date{\normalsize First version: September 2, 2026. This version:
September 2, 2026.}

\begin{document}
\maketitle

\begin{abstract}
\noindent A prompt that fails a handful of scored attempts is rated
impossible by the extrapolation $1-(1-\hat p)^k$, and rated so at
every $k$. The ingredients of the repair are established: the bias
was flagged by \citet{chen2021}, the population-level fix by a
fitted difficulty prior is published \citep{kazdan2025}, and the
per-prompt posterior predictive measured here is precisely the
estimator of \citet{verine2026}, who condition on each task's count
and average the posterior predictions. This note is an independent
evaluation of that estimator, on released rollout records with a
genuinely extrapolative design:
two training attempts per prompt, six held out, nine thousand
prompts. The plug-in understates the aggregate pass@6 by seventeen
points and, used as a reachability screen at the threshold a
published diagnostic uses, declares $5{,}511$ prompts unreachable of
which a third succeed. The posterior predictive removes most of the
damage, and the study separates how: shrinking the point estimate
accounts for the bulk of the log-loss improvement, posterior
integration for a further measurable step, and a fixed Jeffreys
prior captures nearly all of the fitted families' gain. Limits are
measured rather than assumed, including per-count calibration and
the insensitivity of prompt rankings.
\end{abstract}

\section{Introduction}

Language models are improved after pretraining by having them
attempt problems and scoring the attempts automatically: a unit test
passes, a final answer matches, a verifier accepts. The evaluations
that steer this process sample each prompt several times, because a
model that solves a problem once in sixteen tries can often be
trained into solving it reliably. Pass@$k$, the probability that at
least one of $k$ sampled attempts succeeds, averaged over prompts,
is the working currency of post-training evaluation, and decisions
of real consequence ride on it: which prompts to train on, whether a
fine-tuning recipe helped, when to stop sampling.

When $n \ge k$ scored attempts per prompt are available, the
combinatorial estimator of \citet{chen2021},
$1-\binom{n-c}{k}/\binom{n}{k}$ with $c$ the correct count, is
unbiased for the per-prompt coverage, and unbiasedness for $k \le n$
is as much as frequentist estimation can ask. It is not the same as
optimal prediction: a decision scored by log loss or Brier score is
a different objective, and the distinction matters below. The device
this note measures arises when a coverage number is wanted at depth
$k$ from an estimate $\hat p = s/m$: evaluate the formula at the
point estimate and report $1-(1-\hat p)^k$: the plug-in, in the
terminology of \citet{verine2026}, who call it the naive plug-in
estimator. \citet{malladi2026} use it inside a coverage
diagnostic,
writing ``we convert pass@1 to an estimated pass@16 value using
$f_{16}(p):=1-(1-p)^{16}$''; their pass@1 estimates rest on at least
sixteen samples per problem per evaluation seed across three seeds
(their Appendix D.3), so theirs is an instance of the plug-in
transformation at substantial sampling depth rather than
extrapolation beyond it, a distinction this note keeps. The
structural point transfers to their setting; the magnitudes measured
here do not.

\citet{chen2021} already cautioned that the plug-in is biased, and
their combinatorial estimator addresses $k \le n$. The
population-level repair by a fitted difficulty prior has been
published: \citet{kazdan2025} fit a beta prior by marginal maximum
likelihood and report the prior expectation of coverage, with
zero-success problems handled by prior mass rather than declared
impossible. The per-prompt posterior predictive is also published:
\citet{verine2026} write ``we therefore use the posterior
predictive pass@$k$ for each task and average across tasks,'' derive
the beta-function-ratio form of the estimator, and additionally
supply zero-one inflated priors and a theorem ordering population
pass@$k$ at fixed pass@1. Public evaluation code implements the
beta-binomial version as well. This note claims no new estimator.
What it adds is an independent evaluation on new data, under a
design that actually extrapolates and scores per prompt: how large
the plug-in's defects are against held-out batches, how much of the
repair is plain shrinkage versus posterior integration, whether a
fixed prior suffices, where the fitted models are and are not
calibrated, and what changes in a threshold decision of the kind
diagnostics actually take. A probit-normal prior is run alongside
the beta and performs comparably; no advantage is claimed for it.

\section{Related work}

\citet{brown2024} established repeated sampling as an inference-time
scaling axis and model aggregate coverage curves with exponentiated
power laws, computing per-problem pass rates by the combinatorial
estimator. \citet{schaeffer2025} explain those power laws through
the heavy-tailed distribution of per-problem success probabilities,
fitting a three-parameter Kumaraswamy law to per-problem point
estimates and forecasting coverage through the integral that defines
aggregate pass@$k$. \citet{kazdan2025} fit a beta difficulty prior
by marginal maximum likelihood and report aggregate coverage as the
prior expectation of $1-(1-p)^k$, with a budget-allocation
application. \citet{verine2026} supply the estimator this note
evaluates: binomial counts under a shared beta difficulty prior
fitted by marginal maximum likelihood, the posterior updated with
each task's count, the posterior expectation of future pass@$k$
taken per task and averaged. Their beta-function ratio is
algebraically the rising-factorial ratio below, and they distinguish
the prior expectation of coverage from the average of posterior
predictions, the same distinction Proposition~\ref{prop:bayes}'s
cautions require. They go further than this note in two respects,
zero-one inflated priors and an ordering theorem for population
pass@$k$ at fixed pass@1. Elsewhere in the
evaluation literature, \citet{hariri2025} replace pass@$k$ reporting
with Bayesian credible intervals at the model level, and
\citet{miller2024} bring frequentist standard errors to eval
reporting. \citet{malladi2026} supply the motivating diagnostic;
nothing here disputes their method, and the question is confined to
the extrapolation step inside it.

\section{The estimation problem}

The practitioner's situation is this. There are $N$ prompts and a
compute budget that allows $m$ scored attempts at each, so the data
are success counts $s_1, \dots, s_N$ with $s_i \in \{0, \dots, m\}$.
Decisions are then taken about a larger budget $k > m$ that will not
be spent yet, or ever: whether a prompt is worth training on depends
on its chance of being solved within $k$ attempts, and whether a
fine-tuning method helped depends on how many prompts cross that
bar. The problem is to choose the function of $s_i$ that stands in
for the unavailable deeper measurement, and the quality of a choice
is measured the way predictions are: per prompt by log loss and
Brier score against later attempts, in aggregate by the error of the
reported curve.

The model under which the problem has an exact answer is minimal.
Each prompt carries a latent per-attempt success probability $p_i$,
drawn independently from a difficulty distribution $G$ on $[0,1]$;
conditional on $p_i$, attempts are independent
$\mathrm{Bernoulli}(p_i)$ trials, so $s_i \sim \mathrm{Bin}(m,
p_i)$. Two estimands follow. Per prompt, the probability that a
fresh batch of $k$ attempts contains a success,
\begin{equation}
q_k(p_i) \;=\; 1-(1-p_i)^{k},
\end{equation}
to be predicted from $s_i$ alone; and in aggregate, the pass@$k$ the
whole population would exhibit,
\begin{equation}
\pi_k \;=\; \int_0^1 \bigl(1-(1-p)^k\bigr)\,dG(p).
\end{equation}
The conditional independence of attempts given $p_i$ is not the
assumption under attack: attempts are drawn with independent seeds,
and the model grants that. What varies across the field is $p_i$
itself, and the whole question is how to respect that variation when
$s_i$ is a noisy glimpse of it.

\section{The plug-in and its failure mode}

At $s_i=0$ the plug-in predicts zero at every $k$, however small $m$
is, so each eventual success on such a prompt costs unbounded log
loss; the mirror failure occurs at $s_i = m$ for prompts that then
fail everything. The aggregate defect was flagged by
\citet{chen2021}; the following restates it with strictness and the
$k=1$ boundary.

\begin{proposition}[The plug-in is biased downward for every
difficulty]\label{prop:bias}
Fix $m \ge 1$ and $k \ge 2$, and let $s \sim \mathrm{Bin}(m,p)$.
Then
\begin{equation}
\mathbb{E}\Bigl[\,1-\bigl(1-\tfrac{s}{m}\bigr)^{k}\Bigr]
\;\le\; 1-(1-p)^k,
\end{equation}
with equality only at $p \in \{0,1\}$; at $k=1$ equality holds for
all $p$. Consequently the expected aggregate plug-in curve lies
below $\pi_k$ for every $G$ not concentrated on $\{0,1\}$.
\end{proposition}

\begin{proof}
Write $\varphi(x) = (1-x)^k$ on $[0,1]$. Its second derivative is
$\varphi''(x) = k(k-1)(1-x)^{k-2} \ge 0$, strictly positive on
$[0,1)$ when $k \ge 2$, so $\varphi$ is convex, and strictly convex
there. The estimate $\hat p = s/m$ satisfies
$\mathbb{E}\,\hat p = p$. Jensen's inequality gives
$\mathbb{E}\,\varphi(\hat p) \ge \varphi(\mathbb{E}\,\hat p)
= (1-p)^k$, and the inequality is strict whenever $\hat p$ is
nondegenerate, which is exactly $p \in (0,1)$. Negating and adding
one reverses the inequality into the display. At $k=1$, $\varphi$ is
affine and Jensen holds with equality. The aggregate statement
follows by integrating the pointwise inequality against $G$ and
taking expectations over the counts.
\end{proof}

The sign is universal; the magnitude is not monotone in $k$. A
direct calculation at $m=2$, $p=0.9$ gives bias $0.0315$ at $k=3$
and $0.0212$ at $k=4$: extrapolating further can cost less, because
the bias also vanishes as $q_k(p)$ saturates toward one. The
experiments below happen to show bias growing over the displayed
range of $k$; that is a fact about this data's difficulty
distribution, not a consequence of convexity.

\section{The posterior predictive}

Give $p$ a prior with two parameters fitted by marginal maximum
likelihood across prompts, and predict with the posterior
expectation
\begin{equation}
\widehat{\mathrm{pass}}_k^{\,\mathrm{post}}
\;=\; \mathbb{E}\bigl[\,1-(1-p)^k \,\bigm|\, s\,\bigr].
\label{eq:post}
\end{equation}
Under a $\mathrm{Beta}(a,b)$ prior the posterior is
$\mathrm{Beta}(a+s,\; b+m-s)$ and
\begin{equation}
\mathbb{E}\bigl[(1-p)^k \mid s\bigr]
= \prod_{j=0}^{k-1} \frac{b+m-s+j}{a+b+m+j},
\end{equation}
a ratio of rising factorials, so Equation~\eqref{eq:post} is
closed-form. Under a probit-normal prior, $p = \Phi(\theta)$ with
$\theta \sim N(\mu, \tau^2)$, the marginal likelihood and the
posterior moment are one-dimensional Gauss--Hermite quadratures over
$\theta$; this is the family that composes with correlated
extensions, since a factor-correlated portfolio of successes
conditions to exactly this form \citep{cotton2026}.

Because $q_k$ is concave, the posterior predictive sits
\emph{below} the plug-in evaluated at the posterior mean:
$\mathbb{E}[q_k(p) \mid s] \le q_k(\mathbb{E}[p \mid s])$. Comparing
those two predictors therefore separates two effects that are easy
to conflate: shrinking the noisy point estimate toward the
population, and integrating the curvature of $q_k$ over the
posterior. The study below prices them separately.

\begin{proposition}[What the posterior predictive is]
\label{prop:bayes}
Let $Y$ indicate a success among $k$ fresh attempts at the same
prompt, drawn independently of the training attempts given $p$.
Under the model, with prior equal to the true difficulty law $G$:
(i) $\mathbb{E}[Y \mid s] =
\mathbb{E}[\,1-(1-p)^k \mid s\,]$, the exact conditional success
probability; (ii) it minimizes expected Brier score and expected log
loss among all predictors that depend on the data only through $s$;
(iii) its population expectation equals $\pi_k$:
$\mathbb{E}_{s}\bigl[\mathbb{E}[q_k(p) \mid s]\bigr] = \pi_k$ when
the outer expectation is taken under the model's marginal law of
$s$.
\end{proposition}

\begin{proof}
Given $p$, $Y$ is Bernoulli with parameter $q_k(p)$ and is
independent of $s$. The tower property gives
$\mathbb{E}[Y \mid s]
= \mathbb{E}\bigl[\mathbb{E}[Y \mid p, s] \bigm| s\bigr]
= \mathbb{E}[q_k(p) \mid s]$, which is (i). For (ii), the
conditional expectation minimizes mean squared error among
$s$-measurable predictors, which is the Brier statement, and the
true conditional probability minimizes expected log loss because the
excess log loss of any other predictor is the expectation of a
Kullback--Leibler divergence, which is nonnegative. For (iii),
$\mathbb{E}\bigl[\mathbb{E}[q_k(p) \mid s]\bigr]
= \mathbb{E}\,q_k(p) = \pi_k$, the tower property again.
\end{proof}

These are standard facts of Bayesian prediction, recorded to fix
what is and is not claimed: they explain the method, and they are
not a contribution. Two cautions attach to (iii). It is a statement
about population expectations under the true model; the average of
posterior predictions over a particular fitted sample need not equal
the prior expectation of coverage under the fitted prior, so the
population estimator of \citet{kazdan2025} is computed as its own
quantity in the tables below. In the present study the two happen to
nearly coincide, because a two-parameter family fitted to counts on
$m=2$ attempts nearly interpolates the three-cell count histogram;
that coincidence is particular, not general.

\section{Numerical study on released rollouts}

The data are the publicly released Pass8 rollout records of the RLVE
training suite: procedurally generated reasoning tasks, each with an
automatic verifier, of the kind used as reinforcement-learning
environments for language models. They carry nine thousand prompts
across eighteen environment families and eight scored samples per
prompt from the Qwen3-4B-Thinking model.\footnote{Dataset:
\texttt{CL-From-Nothing/RLVE-Qwen3-4B-Thinking-2507-Pass8-Rollouts}
on Hugging Face. Extraction and experiment scripts are committed
under \texttt{research/cavity\_calculus/exp2\_passk/} in the
\texttt{winning} repository.} RLVE shapes rewards, so a positive
reward can denote partial credit: of $23{,}152$ positive rewards,
$2{,}988$ lie strictly between zero and one. Success here means
reward at least one; the released rows carry no separate accuracy
field, and rerunning every table under the loose positive-reward
definition changes the numbers but no conclusion.

The design extrapolates. Samples $0$--$1$ are the training half
($m=2$); samples $2$--$7$ are held out, and the per-prompt target is
a success among those six ($k = 6 > m$), a single Bernoulli outcome
scored by log loss and Brier score. The aggregate reference for
$k \le 6$ is the combinatorial estimator on the six held-out samples
only, disjoint from training, and it is a reference estimate, not
truth. The fitted priors are $\mathrm{Beta}(0.31, 0.79)$ and
probit-normal $\mu = -1.06$, $\tau = 1.53$.

\begin{table}[h]
\centering
\begin{tabular}{lcc}
\toprule
predictor of held-out pass@6 & log loss & Brier \\
\midrule
plug-in $1-(1-s/2)^6$ & 6.065 & 0.248 \\
plug-in at the fitted posterior mean & 0.570 & 0.185 \\
Jeffreys $\mathrm{Beta}(\tfrac12,\tfrac12)$ posterior predictive
  & 0.536 & 0.183 \\
fitted beta posterior predictive & 0.510 & 0.172 \\
fitted probit-normal posterior predictive & 0.509 & 0.172 \\
\bottomrule
\end{tabular}
\caption{Per-prompt scores over nine thousand held-out Bernoulli
outcomes, $m=2$ training attempts, $k=6$. The plug-in's log loss is
computed with predictions clipped to $[10^{-12}, 1-10^{-12}]$;
unclipped it is infinite, and the value is a clipping convention,
not a stable measure: clipping at $10^{-6}$, $10^{-3}$, $10^{-2}$
gives $3.11$, $1.63$, $1.14$. The decomposition is the point of the
table: shrinkage alone (row two) removes most of the damage,
posterior integration improves on row two by a further eleven
percent of log loss and seven percent of Brier, and the fixed
Jeffreys prior, with no fitted hyperparameters, captures nearly all
of the fitted families' gain.}
\label{tab:perprompt}
\end{table}

\begin{table}[h]
\centering
\begin{tabular}{lcccc}
\toprule
aggregate & $k=1$ & $k=2$ & $k=4$ & $k=6$ \\
\midrule
reference estimate (six held-out samples) & 0.280 & 0.386 & 0.498
  & 0.558 \\
plug-in from the training half & 0.282 & 0.335 & 0.374 & 0.384 \\
fitted beta, average of posteriors & 0.282 & 0.388 & 0.491 & 0.546 \\
fitted beta, prior expectation & 0.282 & 0.388 & 0.491 & 0.546 \\
\bottomrule
\end{tabular}
\caption{Aggregate pass@$k$ from two training samples per prompt.
The plug-in is $17.4$ points under the reference at $k=6$; the
posterior rows are within $1.3$ points. The prior-expectation row is
the population estimator of \citet{kazdan2025}; it nearly coincides
with the average of posterior predictions here because the fitted
two-parameter family nearly interpolates the three-cell count
histogram at $m=2$, a coincidence of this design rather than an
identity.}
\label{tab:aggregate}
\end{table}

Calibration by count locates what the fitted model does and does not
get right. For the $5{,}511$ prompts with no training success the
fitted beta predicts a held-out success rate of $0.318$ against an
observed $0.343$, a modest undershoot; at counts one and two the
predictions $0.843$ and $0.982$ sit close to the observed $0.834$
and $0.977$, where the plug-in's $0.984$ and $1.000$ overshoot. One
structural limit is worth stating plainly: with a common prior and
common $m$, every prompt with the same count receives the same
prediction, so the ranking of prompts is exactly the ranking by
count under either method. The posterior changes thresholds and
magnitudes, not order.

That is where the decision demonstration lives. The coverage
diagnostic of \citet{malladi2026} works on a base-reachable set,
prompts whose estimated pass@16 lies in $[0.05, 0.95]$. Applying the
lower edge as a reachability screen on predicted pass@6: the plug-in
declares the $5{,}511$ zero-count prompts unreachable, and
$1{,}890$ of them, a third, succeed within the six held-out
attempts; the posterior predictive prices those same prompts at
$0.318$ and declares none unreachable, in line with their observed
$0.343$ rate. A screen built on the plug-in wrongly discards a third
of what it discards at this sampling depth; a screen built on the
posterior predictive makes no such error here, at the cost of
declaring nothing unreachable, which is what a calibrated
probability near one third requires.

\section{Scope}

The bias mechanism and the repair are established prior work; the
measurements are the contribution, and they are bounded by their
setting. The magnitudes scale with the noise in $\hat p$: two
training attempts is an extreme of the small-sample regime, chosen
to make the extrapolation genuine, and \citet{malladi2026} operate
at sixteen or more samples per problem per seed, where the
noise-driven defects measured here are far smaller. The data cover
one model on one environment suite, with success defined by full
reward in the absence of a separate accuracy field; the loose
definition changes every number and no conclusion. The two fitted
families agree closely here, and agreement at $m=2$ establishes
little about the difficulty distribution: both families place no
atom at $p=0$, both therefore predict eventual success with
probability approaching one for every prompt, and two attempts
cannot identify the tail behavior that arbitrary extrapolation would
need; the zero-one inflated priors of \citet{verine2026} address
exactly this. What the study supports is narrower and useful: at
small sampling depth, shrinkage of any reasonable kind repairs most
of the plug-in's per-prompt damage, posterior integration adds a
measured further step, and threshold decisions of the kind
diagnostics take change materially. Two extensions are named and
left undone: measuring the substitution inside a diagnostic at its
native sampling depth on math and code benchmarks, and developing
correlated, heterogeneous attempt portfolios, for which the
probit-normal family composes with factor structure
\citep{cotton2026} but for which nothing is demonstrated here.

\begin{thebibliography}{9}

\bibitem[Brown et~al.(2024)]{brown2024} B.~Brown, J.~Juravsky,
R.~Ehrlich, R.~Clark, Q.~V. Le, C.~R\'e, and A.~Mirhoseini. Large
language monkeys: scaling inference compute with repeated sampling.
arXiv:2407.21787, 2024.

\bibitem[Chen et~al.(2021)]{chen2021} M.~Chen et~al. Evaluating
large language models trained on code. arXiv:2107.03374, 2021.

\bibitem[Cotton(2026)]{cotton2026} P.~Cotton. Scalable inversion of
contests with correlated performances, including softmax and
multinomial probit. arXiv:2609.01133, 2026.

\bibitem[Hariri et~al.(2025)]{hariri2025} A.~Hariri et~al. Don't
pass@$k$: a Bayesian framework for large language model evaluation.
arXiv:2510.04265, 2025.

\bibitem[Kazdan et~al.(2025)]{kazdan2025} J.~Kazdan, R.~Schaeffer,
et~al. Efficient prediction of pass@$k$ scaling in large language
models. arXiv:2510.05197, 2025.

\bibitem[Malladi et~al.(2026)Malladi, Jelassi, Foster, Ash, and
Krishnamurthy]{malladi2026} S.~Malladi, S.~Jelassi, D.~Foster,
J.~T. Ash, and A.~Krishnamurthy. TailSFT: filtered fine-tuning
improves post-training performance. arXiv:2608.25756, 2026.

\bibitem[Miller(2024)]{miller2024} E.~Miller. Adding error bars to
evals: a statistical approach to language model evaluations.
arXiv:2411.00640, 2024.

\bibitem[Schaeffer et~al.(2025)]{schaeffer2025} R.~Schaeffer et~al.
How do large language monkeys get their power (laws)?
arXiv:2502.17578, 2025.

\bibitem[V\'erine et~al.(2026)V\'erine, Le~Bronnec, Allauzen,
and Negrevergne]{verine2026} A.~V\'erine, F.~Le~Bronnec,
A.~Allauzen, and B.~Negrevergne. Estimating pass@$k$ from fewer
samples with hierarchical Bayesian priors. In \emph{1st Workshop on
Combining Theory and Benchmarks (CTB@ICML 2026)}, 2026.

\end{thebibliography}

\end{document}
